% Data. Sources: refs/docs/train_eval_data.md §1--5; HTML §4. Section layout (overview,
% corpus table, pipeline paragraphs) follows the data sections of the 1.0 report and
% Kimi-Audio, per AGENTS.md rule 6. Sampling weights, mixing probabilities and counterpart
% ratios are deliberately left out (user, 2026-10-31).

\section{Data}
\label{sec:data}

Training has two stages (\Cref{sec:training}).
Stage~1, multi-task supervised fine-tuning, uses plain ASR, hotword ASR, synthetic two- and three-talker target-speaker mixtures and turn-taking samples cut from meeting corpora.
This section describes the data of stage~2, target-speaker robustness training.
Target-speaker samples make up most of the stage-2 data: online mixtures, scene pools, turn-taking samples and real meeting windows, plus samples in which nobody speaks or a sentence is surrounded by babble.
The rest are plain and hotword ASR samples, intended to retain those capabilities.
\Cref{tab:corpora} lists the public corpora; the stage-2 data also include a small amount of domain-specific speech and meeting-room noise floors, which are not listed.

\begin{table}[t]
  \caption{Public corpora used in stage~2.}
  \label{tab:corpora}
  \centering
  \small
  \begin{tabularx}{\linewidth}{@{}>{\raggedright\arraybackslash}p{3.4cm}>{\raggedright\arraybackslash}X@{}}
    \toprule
    Use & Corpora \\
    \midrule
    Target-speaker voices & AISHELL-2~\cite{du2018aishell2}, LibriSpeech~\cite{panayotov2015librispeech}, KeSpeech~\cite{tang2021kespeech}, MagicData, Emilia~\cite{he2024emilia} \\
    Interferers and real conversations & AliMeeting~\cite{yu2022m2met}, AISHELL-4~\cite{fu2021aishell4}, MagicData-RAMC~\cite{yang2022open}, AMI~\cite{carletta2005ami}, CHiME-6~\cite{watanabe2020chime6}, Emilia conversations \\
    Plain ASR & WenetSpeech~\cite{zhang2022wenetspeech}, GigaSpeech~\cite{chen2021gigaspeech}, Common Voice~\cite{ardila2020common}, FLEURS, KeSpeech, AISHELL-2, LibriSpeech, 3D-Speaker~\cite{zheng2023speaker}, Voice-in-the-Wild-2M~\cite{xie2026megaasr} \\
    Noise & MUSAN~\cite{snyder2015musan}, WHAM!~\cite{wichern2019wham}, DNS Challenge noise~\cite{dubey2023icassp}, AudioSet, synthetic babble \\
    Room impulse responses & DNS Challenge, OpenSLR 26 and 28, MIT measured rooms~\cite{traer2016statistics} \\
    \bottomrule
  \end{tabularx}
\end{table}
% TODO: add citations for FLEURS, AudioSet and OpenSLR 26/28 (not yet in refs INDEX).

\subsection{Online mixing}
\label{subsec:data_mixing}

Most target-speaker samples are synthesized on the fly when they are loaded, so a source item yields a different mixture each time it is drawn.
A mixture has up to four interferers and varies in signal-to-interference ratio and overall level; speakers may pass through near-field or far-field room responses, and noise is added on top of a noise floor.
The enrollment is a random 1--20\,s crop and may be degraded by another voice, a low level, narrowband filtering, leading silence, telephone codecs or clipping.
The target may enter or leave midway.
For absent samples, the target is removed from the mixture and the transcript is empty.

\subsection{Scene pools and target-absent counterparts}
\label{subsec:data_scenes}

Random mixing does not reproduce structures that are common in meetings and conversations: two people taking turns, someone else picking up a sentence, the target saying only a word or two, or several people talking at once.
We build scene pools for these structures, following three principles.

\paragraph{Every scene type has target-absent counterparts.}
A counterpart keeps the same scene but enrolls another person, sometimes a similar-sounding one, and its reference transcript is empty.
During development, some scene types, such as turn-taking, were almost always target-present in training.
On those scenes the model still transcribed someone in about a quarter of samples after the enrollment was replaced by that of an unrelated person.
Adding counterparts to these scene types lowered the target-absent false-alarm rate from 4.9\% to 1.9\%, measured at the threshold that fixes the empty-output rate at 5\%.

\paragraph{More sources of target speakers.}
Besides AISHELL-2, Mandarin targets come from KeSpeech, in which the same speaker is recorded once in Mandarin and once in dialect, and from MagicData speakers.
Each new source covers all scene types.

\paragraph{The enrollment never contains the target sentence.}
In corpora where a speaker reads the same text several times, the enrollment could be the target sentence itself; such samples are excluded so that the target cannot be copied from the enrollment.

\subsection{Real meeting windows and label checks}
\label{subsec:data_real}

We cut real multi-talker windows from the training sets of AliMeeting, AISHELL-4, MagicData-RAMC and Emilia, and decide from the utterance-level timeline whether the enrolled speaker talks in each window.
Windows labeled ``enrolled speaker absent'' in which a check found the enrolled speaker talking are excluded.

\subsection{No speech and background babble}
\label{subsec:data_babble}

We add samples with only noise, music or babble and an empty target, intended to keep the model from producing text when nobody speaks.
We also add whole sentences with babble before, after and during them, intended to keep the model from dropping a sentence spoken inside babble.
